\documentclass[10pt]{beamer}
\usefonttheme{professionalfonts}
%\usetheme{Rochester}
%\usetheme{AnnArbor}
\usetheme{Szeged}
\useoutertheme{infolines}
\usecolortheme{dolphin}
\usepackage{graphicx}
% \usepackage[english]{babel}
\usepackage{multicol}
\usepackage{color}
%\usepackage{smartdiagram}


%\definecolor{carmine}{rgb}{0.59, 0.0, 0.09}
\usepackage{amsmath}

\usepackage{graphicx}
\usepackage{tikz}
\usetikzlibrary{shapes.callouts}
\usepackage{multirow}

\usepackage[utf8]{inputenc}

\usepackage{enumerate}
\usepackage[shortlabels]{enumitem}
\usepackage{amsmath,amsthm,amssymb,amsfonts,mathrsfs,latexsym,stmaryrd}
\usepackage{tcolorbox}

\usepackage{color}
\usepackage{amssymb, amsmath, amsbsy}
\usepackage{adjustbox,lipsum}
\usepackage{placeins}

\usepackage{textpos}
\usepackage{wrapfig}
\usepackage[spanish, es-tabla]{babel}

\usepackage{caption}
\setlength{\parindent}{15pt}
\hypersetup{hyperfootnotes=false}

\usepackage{url}
\usepackage{float}
\usepackage{tikz}
\usetikzlibrary{patterns}


\title[ICS-294]{Econometría ICS-294\\ \small{Clase 5: Propiedades de MCO y error de medición}}

\author{Francisco Alfaro Medina}
\institute[USM]
{Universidad Técnica Federico Santa María\\ % Your institution for the title page
	\medskip
	\textit{francisco.alfaro@usm.cl} % Your email address

}
\logo{

\includegraphics[width=3cm]{Figuras/logoUSM.png}
}
\date{}
\newcommand\Fontvi{\fontsize{9}{7.2}\selectfont}
\begin{document}

\begin{frame}
	\titlepage % Print the title page as the first slide
\end{frame}

\begin{frame}{Objetivos de la Clase}

Al finalizar esta clase, el estudiante será capaz de:

\begin{itemize}
    \item[\checkmark] Enunciar los \textbf{supuestos de Gauss-Markov} y entender su rol.
    \item[\checkmark] Derivar la \textbf{varianza} del estimador MCO $\text{Var}(\hat\beta_j)$ y sus errores estándar.
    \item[\checkmark] Demostrar que $s^2$ es un \textbf{estimador insesgado} de $\sigma^2$.
    \item[\checkmark] Probar el \textbf{Teorema de Gauss-Markov}: MCO es BLUE.
    \item[\checkmark] Analizar el efecto del \textbf{error de medición} en variables dependiente e independiente (sesgo de atenuación).
\end{itemize}

\end{frame}

% ===================== SECCION TEORICA: PROPIEDADES DE MCO =====================

\section{Propiedades de MCO: varianza y eficiencia}

\begin{frame}{Los supuestos de Gauss-Markov}

Para el modelo lineal $y_i=\beta_0+\beta_1 x_{i1}+\cdots+\beta_k x_{ik}+u_i$:

\begin{enumerate}
    \item[(SLM.1)] Linealidad en los parámetros: $y=\mathbf{X}\beta+u$.
    \pause
    \item[(SLM.2)] Muestreo aleatorio de $\{(x_i,y_i)\}$.
    \pause
    \item[(SLM.3)] No multicolinealidad perfecta en $\mathbf{X}$ (rango completo).
    \pause
    \item[(SLM.4)] Media condicional cero: $E[u_i \mid \mathbf{X}]=0$ (exogeneidad).
    \pause
    \item[(SLM.5)] Homocedasticidad: $\text{Var}(u_i \mid \mathbf{X})=\sigma^2$.
    \pause
    \item[(SLM.6)] Sin autocorrelación: $\text{Cov}(u_i,u_j\mid \mathbf{X})=0$, $i\neq j$.
\end{enumerate}

\vspace{0.3cm}
\begin{block}{Roles}
SLM.1--SLM.4 garantizan \textbf{insesgamiento}; añadir SLM.5--SLM.6 garantiza \textbf{mínima varianza} (eficiencia).
\end{block}

\end{frame}

\begin{frame}[allowframebreaks]{Derivación: $\text{Var}(\hat\beta_j)$ y errores estándar}

\begin{alertblock}{Resultado: varianza de los estimadores MCO}
Bajo SLM.1--SLM.5, la varianza condicional del estimador de pendiente $j$ en el modelo múltiple es
\[
\text{Var}(\hat\beta_j\mid \mathbf{X})=\frac{\sigma^2}{\text{SCT}_j\,(1-R_j^2)},
\]
donde $\text{SCT}_j=\sum_i(x_{ij}-\bar x_j)^2$ y $R_j^2$ es el $R^2$ de regresar $x_j$ sobre las demás explicativas.
\end{alertblock}

\begin{exampleblock}{Derivación (caso simple, $k=1$)}
\textbf{Paso 1.} El estimador de pendiente puede escribirse como combinación lineal de los $y_i$:
\[
\hat\beta_1=\sum_i w_i\, y_i, \qquad w_i=\frac{x_i-\bar x}{\sum_j (x_j-\bar x)^2}=\frac{x_i-\bar x}{\text{SCT}_x}.
\]
Estos pesos cumplen $\sum_i w_i=0$ y $\sum_i w_i x_i=1$.

\textbf{Paso 2.} Sustituyendo $y_i=\beta_0+\beta_1 x_i+u_i$ y usando las propiedades de los pesos:
\[
\hat\beta_1=\beta_1+\sum_i w_i\, u_i.
\]

\framebreak

\textbf{Paso 3.} Tomando varianza condicional en $\mathbf{X}$ (los $w_i$ son fijos dado $\mathbf{X}$):
\[
\text{Var}(\hat\beta_1\mid \mathbf{X})=\text{Var}\!\Big(\sum_i w_i u_i \;\Big|\;\mathbf{X}\Big)
=\sum_i w_i^2\,\text{Var}(u_i\mid\mathbf{X})+\sum_{i\neq j} w_i w_j\,\text{Cov}(u_i,u_j\mid\mathbf{X}).
\]

\textbf{Paso 4.} Por SLM.6 las covarianzas son $0$, y por SLM.5 $\text{Var}(u_i\mid\mathbf{X})=\sigma^2$:
\[
\text{Var}(\hat\beta_1\mid\mathbf{X})=\sigma^2\sum_i w_i^2
=\sigma^2\,\frac{\sum_i (x_i-\bar x)^2}{\big(\text{SCT}_x\big)^2}
=\frac{\sigma^2}{\text{SCT}_x}.
\]
En el modelo simple $R_1^2=0$, recuperando el caso general. \hfill $\square$
\end{exampleblock}

\begin{block}{Error estándar}
Como $\sigma^2$ es desconocido, lo reemplazamos por su estimador $s^2$:
\[
\text{ee}(\hat\beta_j)=\sqrt{\widehat{\text{Var}}(\hat\beta_j)}=\frac{s}{\sqrt{\text{SCT}_j\,(1-R_j^2)}}.
\]
Mayor $\sigma^2$, menor variación en $x_j$, o mayor colinealidad ($R_j^2\to1$) $\Rightarrow$ estimador menos preciso.
\end{block}

\end{frame}

\begin{frame}[allowframebreaks]{Insesgamiento de $s^2$ como estimador de $\sigma^2$}

\begin{alertblock}{Proposición}
Bajo SLM.1--SLM.5, el estimador
\[
s^2=\frac{1}{n-k-1}\sum_{i=1}^{n}\hat u_i^2=\frac{\text{SCR}}{n-k-1}
\]
es insesgado: $E[s^2\mid\mathbf{X}]=\sigma^2$.
\end{alertblock}

\begin{exampleblock}{Demostración (sketch, caso regresión simple $k=1$)}
\textbf{Paso 1.} Los residuos cumplen $\hat u_i=u_i-(\hat\beta_0-\beta_0)-(\hat\beta_1-\beta_1)x_i$. Se demuestra que
\[
E\!\Big[\sum_i \hat u_i^2 \;\Big|\;\mathbf{X}\Big]=(n-2)\,\sigma^2.
\]

\framebreak

\textbf{Paso 2 (intuición de los grados de libertad).} Los residuos satisfacen \emph{dos} restricciones lineales derivadas de las ecuaciones normales:
\[
\sum_i \hat u_i=0, \qquad \sum_i x_i\,\hat u_i=0.
\]
Por estas dos restricciones, los $n$ residuos solo tienen $n-2$ ``piezas'' de información libres: estimamos $\sigma^2$ con $n-2$ grados de libertad (en general $n-k-1$).

\textbf{Paso 3.} Dividiendo por los grados de libertad:
\[
E[s^2\mid\mathbf{X}]=\frac{1}{n-2}\,E\!\Big[\sum_i\hat u_i^2\Big|\mathbf{X}\Big]=\frac{(n-2)\sigma^2}{n-2}=\sigma^2.
\]
\hfill $\square$
\end{exampleblock}

\begin{block}{¿Por qué dividir por $n-k-1$ y no por $n$?}
Si dividiéramos por $n$ obtendríamos un estimador \emph{sesgado a la baja} de $\sigma^2$: los residuos MCO son ``más pequeños'' que los errores reales porque MCO minimiza $\sum\hat u_i^2$ por construcción.
\end{block}

\end{frame}

\begin{frame}[allowframebreaks]{Teorema de Gauss-Markov: MCO es BLUE}

\begin{alertblock}{Teorema de Gauss-Markov}
Bajo los supuestos SLM.1--SLM.5, el estimador MCO $\hat\beta_j$ es el de \textbf{mínima varianza} dentro de la clase de estimadores \emph{lineales} (en $y$) e \emph{insesgados}. Es decir, MCO es \textbf{BLUE} (Best Linear Unbiased Estimator).
\end{alertblock}

\begin{exampleblock}{Demostración (caso pendiente, regresión simple)}
\textbf{Paso 1 (clase de estimadores).} Consideramos cualquier estimador lineal en $y$:
\[
\tilde\beta_1=\sum_i c_i\, y_i,
\]
con pesos $c_i$ que pueden depender de $\mathbf{X}$ pero no de $y$. Escribimos $c_i=w_i+d_i$, donde $w_i=(x_i-\bar x)/\text{SCT}_x$ son los pesos de MCO y $d_i$ es una desviación arbitraria.

\framebreak

\textbf{Paso 2 (condición de insesgamiento).} Sustituyendo $y_i=\beta_0+\beta_1 x_i+u_i$ y tomando esperanza (con $E[u_i\mid\mathbf{X}]=0$):
\[
E[\tilde\beta_1\mid\mathbf{X}]=\beta_0\sum_i c_i+\beta_1\sum_i c_i x_i.
\]
Para que $\tilde\beta_1$ sea \emph{insesgado para todo} $\beta_0,\beta_1$ se requiere
\[
\sum_i c_i=0 \qquad\text{y}\qquad \sum_i c_i x_i=1.
\]
Como los $w_i$ ya cumplen esto, las desviaciones deben satisfacer
\[
\sum_i d_i=0, \qquad \sum_i d_i x_i=0.
\]

\framebreak

\textbf{Paso 3 (varianza del estimador alternativo).} Bajo homocedasticidad y no autocorrelación (SLM.5--SLM.6):
\[
\text{Var}(\tilde\beta_1\mid\mathbf{X})=\sigma^2\sum_i c_i^2=\sigma^2\sum_i (w_i+d_i)^2
=\sigma^2\Big(\sum_i w_i^2+2\sum_i w_i d_i+\sum_i d_i^2\Big).
\]

\textbf{Paso 4 (el término cruzado se anula).} Usando $w_i=(x_i-\bar x)/\text{SCT}_x$ y las restricciones del Paso 2:
\[
\sum_i w_i d_i=\frac{1}{\text{SCT}_x}\sum_i (x_i-\bar x)\,d_i
=\frac{1}{\text{SCT}_x}\Big(\underbrace{\sum_i x_i d_i}_{=0}-\bar x\underbrace{\sum_i d_i}_{=0}\Big)=0.
\]

\framebreak

\textbf{Paso 5 (conclusión).} Por lo tanto
\[
\text{Var}(\tilde\beta_1\mid\mathbf{X})=\underbrace{\sigma^2\sum_i w_i^2}_{=\,\text{Var}(\hat\beta_1\mid\mathbf{X})}+\;\sigma^2\sum_i d_i^2
\;\ge\; \text{Var}(\hat\beta_1\mid\mathbf{X}),
\]
ya que $\sigma^2\sum_i d_i^2\ge 0$. La igualdad se alcanza solo si $d_i=0$ para todo $i$, es decir, cuando $\tilde\beta_1=\hat\beta_1$.
\\[4pt]
Así, \textbf{ningún} estimador lineal e insesgado tiene menor varianza que MCO. \hfill $\square$
\end{exampleblock}

\begin{block}{Advertencia}
Gauss-Markov \emph{no} requiere normalidad de los errores. Pero si falla SLM.5 (heterocedasticidad) o SLM.6 (autocorrelación), MCO sigue insesgado pero \textbf{ya no es BLUE}: existen estimadores (p.ej.\ MCG) con menor varianza.
\end{block}

\end{frame}

% ===================== ERROR DE MEDICION (CONTENIDO ORIGINAL) =====================

\section{Error de Medición}

\begin{frame}{El mundo real}
    \begin{itemize}
        \item Hasta ahora, hemos hablado de las variables $y$ e $x$ cómo dadas en nuestro modelo.
        \item Pero una vez que tratamos de estimar la relación entre estas variables, hay que usar lo que está disponible en los datos reales.
        \item A menudo, no se cuenta con datos sobre la variable económica que realmente nos interesa.
        \item Ejemplo
        \begin{itemize}
            \item Ingreso verdadero vs. ingreso declarado
            \item Calor\'ias consumidas vs. calor\'ias compradas
        \end{itemize}

    \end{itemize}
\end{frame}

\begin{frame}{Error de Medición}
    \begin{itemize}
        \item La variable que nos interesa es $W^{*}$, pero solo tenemos una versión imprecisa de esta variable donde:
$$
W=W^{*}+u
$$
\item Se define como ``error de medida" $u=W-W^{*}$. Es decir, la diferencia entre la variable que observamos y la verdadera variable.
\item Supongamos el error menos grave posible. Lo llamamos error cl\'asico cuando $E(u)=0$ y $\operatorname{Cov}(u, W^{*})=0$.
\item ¿Qué pasa si tenemos errores de medición cl\'asico en la variable dependiente? ¿E independiente?
    \end{itemize}
\end{frame}

\section{Error de medición en la variable dependiente}

\begin{frame}{Error de medición en la variable dependiente}
    \begin{itemize}
        \item Considere el siguiente modelo verdadero,  no hay sesgo por variable omitida:

        \begin{eqnarray}
                                   Y^{*}=\beta_{0}+\beta_{1} X+\varepsilon\\
                                   E[\varepsilon \mid X]=0  \nonumber
       \end{eqnarray}
        \item Pero en lugar de $Y^{*}$ solamente tenemos $Y=Y^{*}+u$. Donde el error de medici\'on lo asumimos cl\'asico: $\operatorname{Cov}\left(Y^{*}, u\right)=0$ y $V(u)=\sigma_{u}^{2}$.
        \item Si sustituimos $Y-u=Y^{*}$ en (1) y despejamos, podemos estimar:

        \begin{equation}
             Y=\beta_{0}+\beta_{1} X+v        \end{equation}


        \item Donde $v=(\varepsilon+u)$
        \item ¿Ser\'a $\hat{\beta}_{1}$ insesgado?
        \item Como suponemos que el error de medici\'on de la variable $Y$ {\color{blue} no} esta correlacionado con X, entonces no causar\'a sesgo.
    \end{itemize}
\end{frame}

\begin{frame}{Error de medición en la variable dependiente}
    \begin{itemize}
              \item Para que los estimadores $\hat{\beta}_{0}$ y $\hat{\beta}_{1}$ sean insesgados, necesitamos que
        $$
        E\left(\hat{\beta}_{1}\right)=\beta_{1}
        $$
        \item Sabemos que
        $$
        \begin{gathered}
        E\left(\hat{\beta}_{1}\right)=\frac{\operatorname{Cov}(X, Y)}{\operatorname{Var}(X)}=\frac{\operatorname{Cov}\left(X, Y^{*}+u\right)}{\operatorname{Var}(X)} \\
        E\left(\hat{\beta}_{1}\right)=\frac{\operatorname{Cov}\left(X, Y^{*}\right)}{\operatorname{Var}(X)}+\frac{\operatorname{Cov}(X, u)}{\operatorname{Var}(X)} \\
        E\left(\hat{\beta}_{1}\right)=\beta_{1}
        \end{gathered}
        $$
        \item {\color{blue}Dado que $\operatorname{Cov}(X, u)=0$ por la definición del error clásico.}
    \end{itemize}
\end{frame}

\begin{frame}{Intuición}
    Si el error de medida en la variable dependiente no está sistemáticamente relacionado con las variables explicativas, entonces el estimador MCO sigue siendo insesgado.
\end{frame}
\begin{frame}{Impacto sobre precisión}
    \begin{itemize}
        \item  Sin embargo, los estimadores de MCO van a tener mayores varianza


        \item Con errores de medición en la variable dependiente, el t\'ermino del error es $\varepsilon+u$. Y entonces, su varianza es $\operatorname{Var}(\varepsilon)+\operatorname{Var}(u)>\operatorname{Var}(\varepsilon)$.
    \end{itemize}
\end{frame}

\section{Errores de medición en variables independientes}

\begin{frame}{Ejemplo: Relación entre salario e IQ}
    \begin{figure}
        \centering
        \includegraphics[width=0.75\linewidth]{complemetario/Imagenes/RelacionSalarioIQ.jpeg}
    \end{figure}
\end{frame}

\begin{frame}{Errores de medición en variables independientes}
    \begin{itemize}
        \item Cuando la variable mal medida es la independiente, el problema es m\'as grave.
        \item La variable que nos interesa es $X^{*}$, pero solo tenemos una versión imprecisa de esta variable donde
        $$
        X=X^{*}+u
        $$
        \item Vamos a seguir suponiendo errores clásicos, es decir que $\operatorname{Cov}\left(X^{*}, u\right)=0$ y $V(u)=\sigma_{u}^{2}$.
    \end{itemize}
\end{frame}

\begin{frame}{Impacto sobre MCO}
    \begin{itemize}
        \item El modelo verdadero es:
        \begin{equation*}
        Y=\beta_{0}+\beta_{1} X^{*}+\epsilon \tag{1}
        \end{equation*}
        \item Donde asumimos que $E\left[\epsilon \mid X^{*}\right]=0$
        \item Pero podemos estimar:
        \begin{equation*}
            Y=b_{0}+b_{1} X+v \tag{2}
        \end{equation*}

        \item $v=\epsilon-\beta_1u$
        \item Vamos a ver que ahora el estimador $\hat{b}_1$ va a ser sesgado.
    \end{itemize}
\end{frame}

\begin{frame}{Sesgo de atenuaci\'on - sesgo por error de medici\'on en X}
    \begin{itemize}
        \item Estimamos:
        $$
        \begin{gathered}
        Y=b_{0}+b_{1} X+v \\
        E\left[\hat{b}_{1}\right]=\frac{\operatorname{Cov}(X, Y)}{\operatorname{Var}(X)}=\frac{\operatorname{Cov}\left(X^{*}+u, Y\right)}{\operatorname{Var}\left(X^{*}+u\right)}
        \end{gathered}
        $$
        \item Dado que el error es clásico: $\operatorname{Cov}(u, Y)=0$ por lo que:
        $$
        E\left[\hat{b}_{1}\right]=\frac{\operatorname{Cov}\left(X^{*}, Y\right)}{\operatorname{Var}\left(X^{*}+u\right)}
        $$
        \item Además, el denominador lo podemos escribir como la suma de varianzas:
        $$
        \operatorname{Var}\left(X^{*}+u\right)=\operatorname{Var}\left(X^{*}\right)+\operatorname{Var}(u)
        $$
        \item Reemplazando $Y$ por el modelo original (1):
        $$
        E\left[\hat{b}_{1}\right]=\frac{\operatorname{Cov}\left(X^{*}, \beta_{0}+\beta_{1} X^{*}+\epsilon\right)}{\operatorname{Var}\left(X^{*}\right)+\operatorname{Var}(u)}
        $$
    \end{itemize}
\end{frame}

\begin{frame}
    \begin{itemize}
        \item separando las covarianzas:
        $$
        E\left[\hat{b}_{1}\right]=\frac{\beta_{1} \operatorname{Cov}\left(X^{*}, X^{*}\right)+\operatorname{Cov}\left(X^{*}, \epsilon\right)}{\operatorname{Var}\left(X^{*}\right)+\operatorname{Var}(u)}
        $$
        \item Sabemos que el modelo real esta bien especificado por lo que $\operatorname{Cov}\left(X^{*}, \epsilon\right)=0$. Ademas, por definici\'on, $\operatorname{Cov}\left(X^{*}, X^{*}\right)=\operatorname{Var}\left(X^{*}\right)$.
        \item Entonces:
        $$
        E\left(\hat{b}_{1}\right)=\beta_{1} \frac{\operatorname{Var}\left(X^{*}\right)}{\operatorname{Var}\left(X^{*}\right)+\operatorname{Var}(u)}=a \cdot \beta_{1}
        $$
        \item Noten que $a<1$.
        \item Por eso se llama sesgo de atenuaci\'on. El coeficiente de MCO ser\'a m\'as pequeño en valor absoluto, subestimaremos el efecto.
    \end{itemize}
\end{frame}

\begin{frame}{Sesgo de atenuación}
    $$
    E\left(\hat{b}_{1}\right)=\beta_{1} \frac{\operatorname{Var}\left(X^{*}\right)}{\operatorname{Var}\left(X^{*}\right)+\operatorname{Var}(u)}=a \cdot\beta_{1}
    $$
    \begin{itemize}
        \item Si $\beta_{1}>0$, el estimador será menor al par\'ametro.
        \item Si $\beta_{1}<0$, el estimador será mayor al par\'ametro, menor en valor absoluto.
        \item Entonces, tenderemos a subestimar (en valor absoluto) la magnitud de la pendiente.
        \item La magnitud del sesgo depende de cuan grande es $\operatorname{Var}(u)$ comparado a $\operatorname{Var}(X)$. Esto se llama: "signal-to-noise ratio".
    \end{itemize}
\end{frame}

% ===================== EJERCICIOS RESUELTOS =====================

\begin{frame}{Ejercicio Resuelto 1}

\begin{block}{Ejercicio 1 (varianza y error estándar)}
En una regresión simple del gasto en consumo sobre el ingreso (datos CASEN, $n=200$ hogares) se obtuvo:
\[
\text{SCT}_x=\sum_i(x_i-\bar x)^2=5000, \qquad s^2=\frac{\text{SCR}}{n-2}=8.
\]
Calcule (a) $\widehat{\text{Var}}(\hat\beta_1)$ y (b) el error estándar $\text{ee}(\hat\beta_1)$.
\end{block}

\pause

\begin{exampleblock}{Solución}
\textbf{(a)} En el modelo simple $\widehat{\text{Var}}(\hat\beta_1)=s^2/\text{SCT}_x$:
\[
\widehat{\text{Var}}(\hat\beta_1)=\frac{8}{5000}=0.0016.
\]
\textbf{(b)} El error estándar es la raíz:
\[
\text{ee}(\hat\beta_1)=\sqrt{0.0016}=0.04.
\]
Si $\hat\beta_1=0.6$, el estadístico $t=0.6/0.04=15$, altamente significativo. Mayor $\text{SCT}_x$ (más variación en el ingreso) habría dado un estimador aún más preciso.
\end{exampleblock}

\end{frame}

\begin{frame}{Ejercicio Resuelto 2}

\begin{block}{Ejercicio 2 (sesgo de atenuación)}
El retorno verdadero de la habilidad cognitiva sobre el salario es $\beta_1=0.10$. La habilidad se mide con un test ruidoso: $\text{Var}(X^*)=9$ (señal) y $\text{Var}(u)=3$ (ruido del test). Asumiendo error de medición clásico, ¿qué valor esperaríamos de $E[\hat b_1]$ al usar el test medido?
\end{block}

\pause

\begin{exampleblock}{Solución}
El factor de atenuación es
\[
a=\frac{\text{Var}(X^*)}{\text{Var}(X^*)+\text{Var}(u)}=\frac{9}{9+3}=\frac{9}{12}=0.75.
\]
Por lo tanto
\[
E[\hat b_1]=a\cdot\beta_1=0.75\times 0.10=0.075<0.10.
\]
El error de medición \textbf{atenúa} el coeficiente hacia cero: subestimamos el retorno en un $25\%$. Cuanto mayor el ruido relativo, mayor la atenuación.
\end{exampleblock}

\end{frame}

\begin{frame}[fragile]{Ejemplo aplicado en Python: simulación de atenuación}

\begin{exampleblock}{Comprobando el sesgo de atenuación con statsmodels}
\begin{verbatim}
import numpy as np
import statsmodels.api as sm

np.random.seed(294)
n = 5000
xstar = np.random.normal(0, 3, n)     # Var(X*)=9
y = 2 + 0.10*xstar + np.random.normal(0,1,n)
u = np.random.normal(0, np.sqrt(3), n) # ruido del test
x = xstar + u                          # X medido con error

m = sm.OLS(y, sm.add_constant(x)).fit()
print(m.params[1])   # ~0.075, atenuado desde 0.10
\end{verbatim}
\end{exampleblock}

El coeficiente estimado se acerca a $a\cdot\beta_1=0.075$, confirmando la teoría.

\end{frame}

% ===================== EJERCICIOS PROPUESTOS =====================

\begin{frame}{Ejercicios Propuestos}

\begin{block}{Ejercicio 3 (Propuesto -- rutinario)}
En una regresión simple, $s^2=12$ y $\text{SCT}_x=300$. Calcule $\text{ee}(\hat\beta_1)$. Si además el $R^2_j$ en un modelo múltiple fuese $0.64$, ¿en qué factor aumentaría la varianza respecto al caso $R^2_j=0$?
\end{block}

\begin{block}{Ejercicio 4 (Propuesto -- intermedio)}
Demuestre, usando los pesos $w_i$ del estimador MCO, que $E[\hat\beta_1\mid\mathbf{X}]=\beta_1$ (insesgamiento). Indique en qué paso se usa el supuesto SLM.4.
\end{block}

\begin{block}{Ejercicio 5 (Propuesto -- desafiante)}
Un investigador mide el ingreso con error clásico ($\text{Var}(u)=4$) y la señal verdadera tiene $\text{Var}(X^*)=16$. (a) Calcule el factor de atenuación $a$. (b) Si quisiera reducir el sesgo a la mitad, ¿cuánto debería bajar $\text{Var}(u)$? Discuta por qué el problema es más grave que el error de medición en $Y$.
\end{block}

\end{frame}

\begin{frame}{Conclusiones}
    \begin{itemize}
        \item Cuando trabajamos con datos reales, no siempre medimos exactamente lo que queremos en nuestro modelo.
        \item Errores de medición de la variable dependiente no generan sesgo pero si aumentan varianza, dificultando el reconocimiento de relaciones entre variables.
        \item El error de medición en variables independientes genera sesgo de atenuación: siempre vamos a subestimar la magnitud de la relación entre variables.
    \end{itemize}
\end{frame}

\begin{frame}{Resumen}

\begin{block}{Puntos Clave de la Clase 5}
\begin{enumerate}
    \item Los supuestos SLM.1--SLM.4 dan insesgamiento; SLM.5--SLM.6 añaden eficiencia.
    \item $\text{Var}(\hat\beta_j)=\sigma^2/[\text{SCT}_j(1-R_j^2)]$; el error estándar usa $s$ en vez de $\sigma$.
    \item $s^2=\text{SCR}/(n-k-1)$ es insesgado para $\sigma^2$ (corrección por grados de libertad).
    \item \textbf{Gauss-Markov}: bajo SLM.1--SLM.5 MCO es BLUE (mínima varianza entre lineales insesgados).
    \item Error de medición: en $Y$ no sesga (sube varianza); en $X$ genera \textbf{atenuación} hacia cero.
\end{enumerate}
\end{block}

\begin{alertblock}{Para la próxima clase}
Revisar: estimación e inferencia en el modelo de regresión múltiple y pruebas de hipótesis sobre los coeficientes.
\end{alertblock}

\end{frame}

\end{document}
